\section{Reducing all periodic coverings to finitely many patterns}
\label{sec:localization}

We first remove the lattice and its fundamental domain from the subsequent
probability estimate.  A periodic arrangement of bounded center intensity
has a translated ball containing relatively few centers and many points
with controlled residual vectors.  Rounding those centers produces a list
in a fixed finite family.  The argument counts every occurrence of a center,
so it applies to arbitrarily short or skew periods and to coinciding cosets.

Use the fixed parameters $a,D,v,b,r_-,\eta$ from the preceding section,
and let $1\le R\le n^2$.  Define
\[
 C_{\mathrm{loc}}=\frac{1000}{v},\qquad
 w_a(r)=1+n[1-r^2/a^2]_+ \quad (r\ge0),
 \qquad h_n=\frac{\eta}{100n}.
\]
The weight $w_a$ records the radial cost of a center near a target point.
Our finite alphabet and maximum list length are
\begin{equation}\label{eq:local-grid}
 \mathcal Q_n=h_n\mathbb Z^n\cap B_{D+b+5\eta},\qquad
 N_n=\left\lceil\frac{20R}{v}
       \left(\frac{D+b+5\eta}{a}\right)^n\right\rceil.
\end{equation}
An ordered list $P=(p_1,\ldots,p_m)$ over $\mathcal Q_n$ may contain repeated
entries.  Define $G_P\subseteq B_D$ to consist of the points $y$ satisfying
the following three conditions:
\begin{align}
 |y-p_j|&\ge r_--2\eta &&(1\le j\le m),\label{eq:rounded-inner}\\
 \#\{j:|y-p_j|\le b+3\eta\}&\le2C_{\mathrm{loc}}R,
       \label{eq:rounded-count}\\
 \sum_{j:\,|y-p_j|\le b+3\eta}w_a(|y-p_j|)
       &\le2C_{\mathrm{loc}}R.\label{eq:rounded-weight}
\end{align}
The count condition follows from the weight condition because $w_a\ge1$;
we display both to identify the cardinality and cost bounds used later.
All sums count indices, including repeated entries.  Set
\begin{equation}\label{eq:pattern-family}
 \mathcal F_n=
 \left\{P\in\bigcup_{m=0}^{N_n}\mathcal Q_n^m:
              |G_P|\ge\frac14|B_D|\right\}.
\end{equation}
The dependence of $N_n$ and $\mathcal F_n$ on $R$ is suppressed in the
notation.  Empty lists are allowed, with all their sums equal to zero.

\begin{proposition}[Deterministic localization]\label{prop:localization}
For all sufficiently large $n$, uniformly for $1\le R\le n^2$, the family
$\mathcal F_n$ has the following properties.  With
$\lambda_0=\log(1+D/a)$,
\begin{equation}\label{eq:pattern-entropy}
 N_n\le n^{C_1}\exp(\lambda_0n),\qquad
 \log|\mathcal F_n|
       \le\exp\bigl(\lambda_0n+O(\log n)\bigr),
\end{equation}
where $C_1$ and the implicit constant depend only on the fixed parameters.

Let $\mathcal U\subset S^{n-1}$ be any finite set, and put
\[
 K=bB\cap\bigcap_{u\in\mathcal U}\{x:|u\cdot x|\le1\},\qquad
 K^+=(b+3\eta)B\cap
          \bigcap_{u\in\mathcal U}\{x:|u\cdot x|\le1+3\eta\}.
\]
Suppose that $\Lambda$ is any full-rank lattice, $T\subset\mathbb R^n$
is finite, and
\[
 \frac{|T|}{\det\Lambda}\le
       \rho_0:=\frac{R}{2v\omega_na^n}.
\]
If $K+\Lambda+T$ covers $\mathbb R^n$, there is a list
$P=(p_1,\ldots,p_m)\in\mathcal F_n$
such that
\begin{equation}\label{eq:localized-cover}
 B_D\subseteq\bigcup_{j=1}^{m}(p_j+K^+).
\end{equation}
The constructed list in fact satisfies $|G_P|\ge|B_D|/2$.

For every $P\in\mathcal F_n$ and $y\in G_P$, list all residual vectors
$x_i=y-p_j$ with $|y-p_j|\le b+3\eta$, retaining their indices, and denote
their number by $m_y$.  Then
\begin{equation}\label{eq:local-residual-budget}
 \begin{gathered}
 r_--2\eta\le |x_i|\le b+3\eta,\qquad
 m_y\le2C_{\mathrm{loc}}R,\qquad
 \sum_{i=1}^{m_y}w_a(|x_i|)\le2C_{\mathrm{loc}}R,\\
 m_y+\sum_{i=1}^{m_y}w_a(|x_i|)\le4C_{\mathrm{loc}}R.
 \end{gathered}
\end{equation}
\end{proposition}

\begin{proof}
We prove the localization and rounding assertions first, then count the
resulting lists.  All dimension thresholds in the proof are uniform for
$1\le R\le n^2$.

\paragraph{Averaging the full multiset of centers.}
Give each center an index $(t,\ell)\in T\times\Lambda$, at position
$t+\ell$.  Distinct indices remain distinct even if they have the same
position.  This multiset is locally finite because $T$ is finite and
$\Lambda$ is discrete.  Write $\rho=|T|/\det\Lambda$, and let $F$ be a
fundamental parallelepiped of $\Lambda$.  For every nonnegative integrable
function $f$,
\begin{equation}\label{eq:periodic-averaging}
 \frac1{|F|}\int_F
       \sum_{(t,\ell)\in T\times\Lambda} f(y-t-\ell)\,dy
       =\rho\int_{\mathbb R^n}f(x)\,dx.
\end{equation}
Indeed, Tonelli's theorem permits the sums to be interchanged with the
integral.  For each $t$, the sets $F-t-\ell$, $\ell\in\Lambda$, partition
$\mathbb R^n$ up to null sets.  This proves the identity without any
restriction on the shape or size of $F$.

Put $\beta=b+4\eta$.  Call a point $y\in\mathbb R^n$ unrounded-good if
there is no center occurrence at distance less than $r_-$, at most
$C_{\mathrm{loc}}R$ occurrences have distance at most $\beta$, and
\begin{equation}\label{eq:unrounded-weight}
 \sum_{(t,\ell):\,|y-t-\ell|\le\beta}
         w_a(|y-t-\ell|)\le C_{\mathrm{loc}}R.
\end{equation}
Let $G$ be this periodic measurable set.  These conditions respectively
exclude short residuals, bound the number of relevant residuals, and bound
their total radial cost.  Using the full ball ensures that every index
entering a rounded test is counted before rounding.

For a uniform point modulo $\Lambda$, the expected number of occurrences
in its open $r_-$-ball is at most
\[
 \rho\omega_n r_-^n
 \le\frac{R}{2v}
       \left(1-\frac{10\log n}{n}\right)^n
 \le\frac{R n^{-10}}{2v}
 \le\frac1{200}
\]
for large $n$.  Hence the first condition fails with probability at most
$1/200$.  Also
\[
 \left(\frac\beta a\right)^n
 \le\exp(1+4/n^3)<3
\]
for large $n$, so the expected outer count is at most $3R/(2v)$.
For the weighted condition, polar integration gives the exact identity
\begin{align}
 \int_{B_\beta}w_a(|x|)\,dx
 &=\omega_n\beta^n
      +n\int_{B_a}(1-|x|^2/a^2)\,dx\notag\\
 &=\omega_n\beta^n+\frac{2n}{n+2}\omega_na^n
 <5\omega_na^n.\label{eq:weighted-ball-integral}
\end{align}
Equation \eqref{eq:periodic-averaging} therefore bounds the expected
weighted sum by $5R/(2v)$.  Markov's inequality bounds the two remaining
failure probabilities by $3/2000$ and $5/2000$, respectively.  Consequently
\begin{equation}\label{eq:periodic-good-volume}
 \frac{|G\cap F|}{|F|}
 \ge1-\frac1{200}-\frac3{2000}-\frac5{2000}
 >\frac{99}{100}.
\end{equation}

\paragraph{One window with both required bounds.}
For $z$ uniform modulo $\Lambda$, let $M(z)$ count all center occurrences
in $z+B_{D+b+4\eta}$.  Equations \eqref{eq:periodic-averaging} and
\eqref{eq:local-grid} imply
\[
 \mathbb{E} M(z)
 \le\frac{R}{2v}\left(\frac{D+b+4\eta}{a}\right)^n
 \le\frac{N_n}{40}.
\]
Moreover, periodicity and Fubini's theorem give
\[
 \mathbb{E}\frac{|(z+B_D)\setminus G|}{|B_D|}
 =1-\frac{|G\cap F|}{|F|}\le\frac1{100}.
\]
Both expressions concern the entire Euclidean window, including all its
overlapping period copies.  Markov's inequality gives
\[
 \mathbb{P}\{M(z)>N_n\}\le\frac1{40},\qquad
 \mathbb{P}\{|(z+B_D)\setminus G|>|B_D|/2\}\le\frac1{50}.
\]
The sum is less than one.  Fix a slide $z$ outside these two exceptional
events, and retain every occurrence in $z+B_{D+b+4\eta}$.  There are at
most $N_n$ such occurrences, and at least half of $z+B_D$ lies in $G$.

\paragraph{Rounding and transferring coverage.}
Translate this window and the retained centers by $-z$, and round every
center coordinatewise to the nearest point of $h_n\mathbb Z^n$, using a
fixed rule for ties.  Each center moves a distance at most
\begin{equation}\label{eq:rounding-error}
 \delta_n=\frac{\sqrt n h_n}{2}
         =\frac{\eta}{200\sqrt n}\le\frac\eta{200}.
\end{equation}
The rounded centers belong to $\mathcal Q_n$; enumerate them as an ordered
list $P$, retaining all repeated entries.

We check that every $y\in B_D$ with $y+z\in G$ belongs to $G_P$.
Every unrounded residual has length at least $r_-$, so every rounded one
has length at least $r_--\delta_n\ge r_--2\eta$.
A rounded residual of length at most $b+3\eta$ comes from an unrounded
residual of length at most
\[
 b+3\eta+\delta_n<b+4\eta.
\]
Thus at most $C_{\mathrm{loc}}R$ rounded indices are relevant.  This
strict separation of the two outer radii ensures that every index entering
a rounded test was counted before rounding.

The function $w_a$ is globally Lipschitz on $[0,\infty)$ with constant
$2n/a$.  Its derivative has absolute value $2nr/a^2\le2n/a$ for
$0<r<a$, and it is constant for $r>a$.  If $r$ and $r'$ are the
unrounded and rounded lengths of a relevant residual, then
\[
 w_a(r')\le w_a(r)+\frac{2n\delta_n}{a}
          \le\left(1+\frac{2n\delta_n}{a}\right)w_a(r),
\]
where the last step uses $w_a(r)\ge1$.  Summing over the relevant indices,
all of which were included in \eqref{eq:unrounded-weight}, gives
\[
 \sum_{j:\,|y-p_j|\le b+3\eta}w_a(|y-p_j|)
 \le\left(1+\frac{2n\delta_n}{a}\right)C_{\mathrm{loc}}R
 \le2C_{\mathrm{loc}}R.
\]
This verifies \eqref{eq:rounded-inner}--\eqref{eq:rounded-weight}.
In particular, $|G_P|\ge|B_D|/2$ and $P\in\mathcal F_n$.

Now assume that the original periodic arrangement covers by $K$.
For any point of $z+B_D$, every center whose translate of $K$ contains
that point lies in $z+B_{D+b}$ because $K\subseteq bB$.  Such a center was
retained.  Moving it by at most $\delta_n$ changes the residual norm by at
most $\delta_n$ and each scalar product with a unit normal by at most
$\delta_n$.  Hence its rounded translate of $K^+$ contains the translated
point of $B_D$.  This proves \eqref{eq:localized-cover} for every point of
the closed window.

\paragraph{Counting the lists.}
Write $J_n=|\mathcal Q_n|$.  A containing cube gives
\[
 J_n\le\left(1+\frac{2(D+b+5\eta)}{h_n}\right)^n,
 \qquad \log J_n=O(n\log n).
\]
Since $b=a(1+1/n)$, $\eta=n^{-4}$, and $1\le R\le n^2$,
\[
 \log N_n=\lambda_0n+O(\log n).
\]
This proves the first bound in \eqref{eq:pattern-entropy}.  Counting ordered
lists allows all multiplicities and gives
\[
 |\mathcal F_n|\le\sum_{m=0}^{N_n}J_n^m
       \le(N_n+1)J_n^{N_n}.
\]
The two preceding estimates give the second bound in
\eqref{eq:pattern-entropy}.
Finally, \eqref{eq:local-residual-budget} follows directly from the
definition of $G_P$ and $w_a\ge1$.
\end{proof}

The family $\mathcal F_n$ depends only on the fixed numerical parameters,
$n$, and $R$.  Proposition~\ref{prop:localization} holds for every choice
of the slab normals and every periodic covering of the required intensity.
Thus a probability bound uniform over this finite family also controls
centers chosen after the body is known.  The next sections obtain such a
bound by finding a point of $G_P$ missed by all translates in a fixed list.
